
\documentclass[aspectratio=169, handout, t]{beamer}

%\usepackage[table]{xcolor}
\mode<presentation> {
\setbeamercovered{transparent}
  \usetheme{Boadilla}

%  \usetheme{Pittsburgh}
%\usefonttheme[2]{sans}
\renewcommand{\familydefault}{cmss}
%\usepackage{lmodern}
%\usepackage[T1]{fontenc}
%\usepackage{palatino}
%\usepackage{cmbright}
\usepackage{bm}
\usepackage{listings}
\useinnertheme{rectangles}
}
\usepackage{amsmath}
\providecommand{\Var}{\mathrm{Var}}
\usepackage{tcolorbox}
\setbeamercolor{normal text}{fg=black}
\setbeamercolor{structure}{fg= blue}
\definecolor{trial}{cmyk}{1,0,0, 0}
\definecolor{trial2}{cmyk}{0.00,0,1, 0}
\definecolor{darkgreen}{rgb}{0,.4, 0.1}
\usepackage{array}
\beamertemplatesolidbackgroundcolor{white}  \setbeamercolor{alerted
text}{fg=red}
\setbeamertemplate{caption}[numbered]\newcounter{mylastframe}
\usepackage{booktabs}

\lstset{%
  language=R,
  basicstyle=\ttfamily\small,
  keywordstyle=\color{blue},
  commentstyle=\color{darkgreen},
  stringstyle=\color{red},
  showstringspaces=false,
  breaklines=true,
  frame=single,
  backgroundcolor=\color{gray!10}
}

\font\domino=domino
\def\die#1{{\domino#1}}
\usepackage{tikz}
\usetikzlibrary{arrows}
\usepackage{colortbl}

\renewcommand{\familydefault}{cmss}
%\usepackage[all]{xy}

\usetikzlibrary{decorations.pathreplacing}
\usepackage{lipsum}


 \newenvironment{changemargin}[3]{%
 \begin{list}{}{%
 \setlength{\topsep}{0pt}%
 \setlength{\leftmargin}{#1}%
 \setlength{\rightmargin}{#2}%
 \setlength{\topmargin}{#3}%
 \setlength{\listparindent}{\parindent}%
 \setlength{\itemindent}{\parindent}%
 \setlength{\parsep}{\parskip}%
 }%
\item[]}{\end{list}}
\usetikzlibrary{arrows}

\usecolortheme{lily}

\newtheorem{com}{Comment}
\newtheorem{lem} {Lemma}
\newtheorem{prop}{Proposition}
\newtheorem{condition}{Condition}
\newtheorem{thm}{Theorem}
\newtheorem{defn}{Definition}
\newtheorem{cor}{Corollary}
\newtheorem{obs}{Observation}
 \numberwithin{equation}{section}

\makeatletter
\def\beamerorig@set@color{%
  \pdfliteral{\current@color}%
  \aftergroup\reset@color
}
\def\beamerorig@reset@color{\pdfliteral{\current@color}}
\makeatother
\setbeamertemplate{navigation symbols}{}

\useoutertheme{miniframes}
\title[PLSC 30700]{Linear Models Lecture 11: Asymptotic theory for OLS}

\author{Robert Gulotty}
\institute[Chicago]{University of Chicago}
\vspace{0.3in}


\begin{document}

\begin{frame}
\maketitle
\end{frame}

%%%%%%%%%%%%%%%%%%%%%%%%%%%%%%%%%%%%%%%%%%%%%%%%%%%%%%%%%%%%
\section{Finishing the CLT: Moment Generating Functions}
\setcounter{subsection}{1}
%%%%%%%%%%%%%%%%%%%%%%%%%%%%%%%%%%%%%%%%%%%%%%%%%%%%%%%%%%%%

\begin{frame}{Where We Left Off}

\textbf{Last lecture.} Stated the CLT:
\[
\sqrt n\,(\bar X_n - \mu) \;\overset{d}{\to}\; N(0,\sigma^2).
\]
True for any iid sample with finite variance; we did not prove it.

\pause

\medskip

\textbf{Today.} The proof requires one tool we have not developed: the \textbf{moment generating function}. We will:
\begin{enumerate}\setlength\itemsep{2pt}
\item Define the MGF and work a complete example.
\item Sketch the CLT proof in three steps.
\item Turn to the delta method and OLS asymptotics.
\end{enumerate}

\end{frame}

%%%%%%%%%%%%%%%%%%%%%%%%%%%%%%%%%%%%%%%%%%%%%%%%%%%%%%%%%%%%

\begin{frame}{What Is a Moment Generating Function?}

\begin{defn}[Moment generating function]
For a random variable $X$,
\[
M_X(t) \;\equiv\; \mathbb{E}\!\bigl[e^{tX}\bigr] \;=\; \int e^{tx}\,f(x)\,dx,
\]
defined for $t$ in some open interval around $0$ where the expectation is finite.
\end{defn}

\pause

\medskip

\textbf{Two facts make the MGF useful.}
\begin{enumerate}\setlength\itemsep{6pt}
\item Its derivatives at $0$ are the moments: $M_X^{(k)}(0) = \mathbb{E}[X^k]$ (next slide).
\item For independent $X$ and $Y$, $M_{X+Y}(t) = M_X(t)\,M_Y(t)$.
\end{enumerate}

\end{frame}

%%%%%%%%%%%%%%%%%%%%%%%%%%%%%%%%%%%%%%%%%%%%%%%%%%%%%%%%%%%%

\begin{frame}{Example 1: Uniform}
If $X\sim U(0,1)$, and $Y\sim U(0,1)$, what is the distribution of $X+Y$?

\end{frame}

\begin{frame}{Example 2: The Distribution of a Sum (Two Fair Dice)}

\textbf{Setup.} Independent rolls $X, Y$, each uniform on $\{1, 2, \ldots, 6\}$. What is the pmf of $X + Y$?

\pause

\medskip

\textbf{Count pairs.} Each of the $36$ pairs $(x, y)$ is equally likely. For each possible sum $s$, count the pairs with $x + y = s$:
\[
\begin{array}{c|ccccccccccc}
s & 2 & 3 & 4 & 5 & 6 & 7 & 8 & 9 & 10 & 11 & 12 \\\hline
\#\text{pairs} & 1 & 2 & 3 & 4 & 5 & 6 & 5 & 4 & 3 & 2 & 1 \\
\end{array}
\]
So $p_{X+Y}(s) = \#\text{pairs}/36$. Triangular, peaked at $7$.

\pause

\medskip

\textbf{The formula.} The count at $s$ is the number of ways $X = x$ can pair with $Y = s - x$, weighted by their joint pmf:
\[
p_{X+Y}(s) \;=\; \sum_x\, p_X(x)\,p_Y(s - x).
\]
This is the \textbf{convolution} of the two pmfs. For continuous $X, Y$, replace the sum with an integral.

\end{frame}

%%%%%%%%%%%%%%%%%%%%%%%%%%%%%%%%%%%%%%%%%%%%%%%%%%%%%%%%%%%%

\begin{frame}{Why the Product Property Matters}

For continuous $X, Y$, the convolution is an integral with the same structure:
\[
f_{X+Y}(s) \;=\; \int f_X(x)\,f_Y(s - x)\,dx.
\]

\pause

\medskip

This integral is awkward to compute once, and worse to iterate $n$ times for a sum $X_1 + \cdots + X_n$.

\pause

\medskip

In MGF terms, the same computation is just a product:
\[
M_{X_1 + \cdots + X_n}(t) \;=\; M_{X_1}(t)\cdots M_{X_n}(t).
\]

\end{frame}

%%%%%%%%%%%%%%%%%%%%%%%%%%%%%%%%%%%%%%%%%%%%%%%%%%%%%%%%%%%%

\begin{frame}{Why ``Moment Generating''? (1 of 2)}

\textbf{Taylor-expand $e^{tX}$ in powers of $t$:}
\[
e^{tX} \;=\; 1 + tX + \frac{(tX)^2}{2!} + \frac{(tX)^3}{3!} + \cdots
\]

\pause

\medskip

\textbf{Take expectations term by term} (valid whenever the MGF exists near $0$):
\[
M_X(t) \;=\; \mathbb{E}[e^{tX}] \;=\; 1 + t\,\mathbb{E}[X] + \frac{t^2}{2!}\mathbb{E}[X^2] + \frac{t^3}{3!}\mathbb{E}[X^3] + \cdots
\]

\pause

\medskip

This is a power series in $t$ whose coefficients are the moments $\mathbb{E}[X^k]/k!$.

\end{frame}

%%%%%%%%%%%%%%%%%%%%%%%%%%%%%%%%%%%%%%%%%%%%%%%%%%%%%%%%%%%%

\begin{frame}{Why ``Moment Generating''? (2 of 2)}

\textbf{Compare to the ordinary Taylor series of $M_X$ at $0$:}
\[
M_X(t) \;=\; M_X(0) + M_X'(0)\,t + \frac{M_X''(0)}{2!}\,t^2 + \frac{M_X'''(0)}{3!}\,t^3 + \cdots
\]

\pause

\medskip

Two power series in $t$ for the same function must have equal coefficients. Match them:
\[
\boxed{\;M_X^{(k)}(0) \;=\; \mathbb{E}[X^k]\; \text{ for every } k \geq 0.\;}
\]

\pause

\medskip

Differentiating $M_X$ at $0$ pulls out the moment of order $k$ --- that is why ``moment generating.'' One analytic function encodes the entire sequence of moments.

\end{frame}

%%%%%%%%%%%%%%%%%%%%%%%%%%%%%%%%%%%%%%%%%%%%%%%%%%%%%%%%%%%%

\begin{frame}{Worked Example: Exponential($\lambda$)}

\textbf{Setup.} $X \sim \text{Exp}(\lambda)$, density $f(x) = \lambda e^{-\lambda x}$ for $x \geq 0$. We know $\mathbb{E}[X] = 1/\lambda$. \emph{Recover it from the MGF.}

\pause

\medskip

\textbf{Step 1 --- Compute the MGF.}
\[
M_X(t) \;=\; \int_0^\infty e^{tx}\,\lambda e^{-\lambda x}\,dx
\;=\; \lambda\int_0^\infty e^{-(\lambda-t)x}\,dx
\;=\; \frac{\lambda}{\lambda - t},\qquad t < \lambda.
\]

\pause

\medskip

\textbf{Step 2 --- Differentiate at $0$.}
\[
M_X'(t) \;=\; \frac{\lambda}{(\lambda-t)^2}
\qquad\Longrightarrow\qquad
M_X'(0) \;=\; \frac{1}{\lambda} \;=\; \mathbb{E}[X]. \;\;\checkmark
\]

\pause

\medskip

No integration of $x\,f(x)$: we computed the MGF once and read the mean off the derivative at $0$. Higher moments work the same way.

\end{frame}

%%%%%%%%%%%%%%%%%%%%%%%%%%%%%%%%%%%%%%%%%%%%%%%%%%%%%%%%%%%%

\begin{frame}{Worked Example: Standard Normal}

\textbf{Setup.} $Z \sim N(0,1)$ has MGF (derive once and commit to memory)
\[
M_Z(t) \;=\; \mathbb{E}[e^{tZ}] \;=\; e^{t^2/2}.
\]
\emph{This is the target of the CLT proof.} Every step of the proof tries to show $M_{Z_n}(t) \to e^{t^2/2}$.

\pause

\medskip

\textbf{Pull off moments by differentiation.}
\begin{align*}
M_Z'(t) &= t\,e^{t^2/2}, & M_Z'(0) &= 0 = \mathbb{E}[Z]. \\
M_Z''(t) &= (1 + t^2)\,e^{t^2/2}, & M_Z''(0) &= 1 = \mathbb{E}[Z^2]. \\
M_Z'''(t) &= (3t + t^3)\,e^{t^2/2}, & M_Z'''(0) &= 0 = \mathbb{E}[Z^3]. \\
M_Z^{(4)}(t) &= (3 + 6t^2 + t^4)\,e^{t^2/2}, & M_Z^{(4)}(0) &= 3 = \mathbb{E}[Z^4].
\end{align*}

\end{frame}

%%%%%%%%%%%%%%%%%%%%%%%%%%%%%%%%%%%%%%%%%%%%%%%%%%%%%%%%%%%%

\begin{frame}{Why MGFs Are the Right Tool for Sums}

\textbf{The key property.} If $X \perp Y$, then
\[
M_{X+Y}(t) = \mathbb{E}[e^{t(X+Y)}] = \mathbb{E}[e^{tX}\,e^{tY}]
\stackrel{\perp}{=} \mathbb{E}[e^{tX}]\,\mathbb{E}[e^{tY}]
= M_X(t)\,M_Y(t).
\]
For an iid sum $S_n = X_1 + \cdots + X_n$: \quad $M_{S_n}(t) = \bigl[M_X(t)\bigr]^n.$

\pause

\medskip

\textbf{Compare the density domain.} Adding independent RVs convolves densities, $f_{X+Y}(z) = \int f_X(x)\,f_Y(z-x)\,dx$.

\pause

\medskip

\emph{This is why the CLT proof uses MGFs:} sums $\to$ products $\to$ powers $\to$ something we can Taylor-expand.

\end{frame}

%%%%%%%%%%%%%%%%%%%%%%%%%%%%%%%%%%%%%%%%%%%%%%%%%%%%%%%%%%%%

\begin{frame}{The MGF Continuity Theorem}

\begin{thm}[L\'evy continuity, MGF version]
Suppose $M_{X_n}(t) \to M_X(t)$ for every $t$ in some open interval around $0$, and $M_X$ is the MGF of a proper distribution. Then $X_n \overset{d}{\to} X.$
\end{thm}

\pause

\medskip

\textbf{Why we need this.} The theorem lets us work at the level of MGFs and \emph{conclude} convergence in distribution for free.

\pause

\medskip

\textbf{The plan for the CLT proof.}
\begin{enumerate}\setlength\itemsep{1pt}
\item Standardize the sample mean as $Z_n$; compute $M_{Z_n}(t)$ via the factorization property.
\item Show $M_{Z_n}(t) \to e^{t^2/2}$.
\item Continuity theorem $\Rightarrow Z_n \overset{d}{\to} N(0,1)$.
\end{enumerate}

\end{frame}

%%%%%%%%%%%%%%%%%%%%%%%%%%%%%%%%%%%%%%%%%%%%%%%%%%%%%%%%%%%%

\begin{frame}{Sketch of the CLT Proof: Setup}

\textbf{Goal.} Show $M_{Z_n}(t) \to e^{t^2/2}$ for $Z_n = \dfrac{\sqrt n\,(\bar X_n - \mu)}{\sigma}$. The MGF continuity theorem then gives $Z_n \overset{d}{\to} N(0,1)$.

\pause

\medskip

\textbf{Standardize.} Let $W_i = (X_i - \mu)/\sigma$. Then
\[
\mathbb{E}[W_i] = 0, \qquad \mathbb{E}[W_i^2] = 1, \qquad Z_n = \tfrac{1}{\sqrt n}\sum_{i=1}^n W_i.
\]

\pause

\medskip

\textbf{Factor via independence.} Since the $W_i$ are iid, the MGF of the sum factors (previous slide):
\[
M_{Z_n}(t) \;=\; \mathbb{E}\!\left[\exp\!\Big(\tfrac{t}{\sqrt n}\sum_i W_i\Big)\right] \;=\; \bigl[M_W(t/\sqrt n)\bigr]^n.
\]

\end{frame}

%%%%%%%%%%%%%%%%%%%%%%%%%%%%%%%%%%%%%%%%%%%%%%%%%%%%%%%%%%%%

\begin{frame}{Sketch of the CLT Proof: Taylor-Expand the MGF}

Taylor-expand $M_W(s)$ around $s=0$:
\[
M_W(s) \;=\; M_W(0) + M_W'(0)\,s + \tfrac{1}{2}M_W''(0)\,s^2 + O(s^3).
\]

\pause

\medskip

Evaluate the derivatives using the moment identities $M_W^{(k)}(0) = \mathbb{E}[W^k]$ (two slides back):
\[
M_W(0) = 1, \qquad M_W'(0) = \mathbb{E}[W] = 0, \qquad M_W''(0) = \mathbb{E}[W^2] = 1.
\]
Hence $M_W(s) = 1 + \tfrac{1}{2}s^2 + O(s^3).$

\pause

\medskip

Plug in $s = t/\sqrt n$:
\[
M_W\!\bigl(t/\sqrt n\bigr) \;=\; 1 + \frac{t^2}{2n} + O(n^{-3/2}).
\]

\end{frame}

%%%%%%%%%%%%%%%%%%%%%%%%%%%%%%%%%%%%%%%%%%%%%%%%%%%%%%%%%%%%

\begin{frame}{Sketch of the CLT Proof: Exponential Limit}

Raise to the $n$-th power (independence):
\[
M_{Z_n}(t) \;=\; \bigl[M_W(t/\sqrt n)\bigr]^n \;=\; \left[1 + \frac{t^2/2}{n} + O(n^{-3/2})\right]^n.
\]

\pause

\medskip

Apply the exponential limit $(1 + x/n)^n \to e^x$. The $O(n^{-3/2})$ term vanishes faster than $1/n$, so it drops out:
\[
M_{Z_n}(t) \;\longrightarrow\; e^{t^2/2}. \qquad \blacksquare
\]

\pause

\medskip

\begin{tcolorbox}[colback=blue!5, colframe=blue!50]
\textbf{The whole proof on one line.} Taylor-expand the MGF, raise to the $n$-th power via independence, take the exponential limit, and invoke continuity.
\end{tcolorbox}

\end{frame}

%%%%%%%%%%%%%%%%%%%%%%%%%%%%%%%%%%%%%%%%%%%%%%%%%%%%%%%%%%%%
\section{The Delta Method}
\setcounter{subsection}{1}
%%%%%%%%%%%%%%%%%%%%%%%%%%%%%%%%%%%%%%%%%%%%%%%%%%%%%%%%%%%%

\begin{frame}{Today's Claims}
\begin{itemize}
\item The least squares estimator $\hat{\bm{\beta}}$ is \emph{consistent} for the projection coefficient $\bm{\beta}$.
\item There is an asymptotic normal approximation to the distribution of $\hat{\bm{\beta}}$ (once normalized).
\item Similarly, the error variance estimators $\hat{\sigma}^2$ and $s^2$ are consistent for $\sigma^2$.
\item There is a consistent estimator of $\bm{V}_\beta$, either under homoskedasticity $\hat{\bm{V}}^0_\beta$, or otherwise.
\item We can use z-statistics for asymptotic standard errors, confidence intervals, and prediction intervals.
\end{itemize}

\pause

\begin{tcolorbox}[colback=blue!5, colframe=blue!50]
\textbf{Strategy:} Apply the tools from Lecture 10 --- WLLN, CLT, CMT --- plus the delta method (introduced today) to derive the asymptotic properties of OLS.
\end{tcolorbox}
\end{frame}

%%%%%%%%%%%%%%%%%%%%%%%%%%%%%%%%%%%%%%%%%%%%%%%%%%%%%%%%%%%%

\begin{frame}{Motivation: Distribution of a Function of Estimators}
\begin{itemize}
\item Suppose we want to estimate some function of parameters $f(b_1, b_2)=b_1/b_2$.\pause
\item What is the distribution of $\bm{f(b)}$?\pause
\item We can use the \textbf{Delta Method} (Asymptotic Normality + Taylor Expansion) to approximate it.\pause
\item The idea: if $\bm{b}$ is approximately normal (by the CLT), then a smooth function $\bm{f(b)}$ is also approximately normal, with variance determined by the derivatives of $f$.
\end{itemize}
\end{frame}

%%%%%%%%%%%%%%%%%%%%%%%%%%%%%%%%%%%%%%%%%%%%%%%%%%%%%%%%%%%%

\begin{frame}{Asymptotic Normality of $\hat{\bm{\beta}}$ (Preview)}
\textbf{Claim} (proved later this lecture): if errors are iid with mean zero and finite variance and $\bm{X}$ is full rank, the CLT gives:
$$\hat{\bm{\beta}}\overset{a}{\sim} N\!\left[\bm{\beta},\ \frac{\sigma^2}{n}  \left[\text{plim}\frac{\bm{X'X}}{n}\right]^{-1}\right], \qquad \text{estimated by } s^2(\bm{X}'\bm{X})^{-1}$$

\pause

\begin{tcolorbox}[colback=blue!5, colframe=blue!50]
For now, take this as given. Once $\hat{\bm{\beta}}$ is approximately normal, we can ask: what is the distribution of a \emph{function} of $\hat{\bm{\beta}}$?
\end{tcolorbox}
\end{frame}

%%%%%%%%%%%%%%%%%%%%%%%%%%%%%%%%%%%%%%%%%%%%%%%%%%%%%%%%%%%%

\begin{frame}{The Delta Method: Linearization}

\textbf{The Jacobian.} For a smooth $\bm{f}: \mathbb{R}^K \to \mathbb{R}^J$, define the $J\times K$ matrix of partial derivatives
\[
\bm{C}(\bm{b}) \;=\; \frac{\partial \bm{f}(\hat{\bm{\beta}})}{\partial \hat{\bm{\beta}}'}.
\]

\pause

\medskip

\textbf{Evaluate at the truth.} Consistency of $\hat{\bm{\beta}}$ plus continuity of $\bm{C}$:
\[
\operatorname{plim}\bm{C}(\hat{\bm{\beta}}) \;=\; \frac{\partial \bm{f}(\bm{\beta})}{\partial \bm{\beta}'} \;\equiv\; \bm{\Gamma}. \qquad\text{(Slutsky)}
\]

\pause

\medskip

\textbf{First-order Taylor expansion} from $\bm{\beta}$ to $\hat{\bm{\beta}}$:
\[
\bm{f}(\hat{\bm{\beta}}) \;=\; \bm{f}(\bm{\beta}) \;+\; \bm{\Gamma}\,(\hat{\bm{\beta}} - \bm{\beta}) \;+\; O_p(n^{-1}).
\]

\end{frame}

%%%%%%%%%%%%%%%%%%%%%%%%%%%%%%%%%%%%%%%%%%%%%%%%%%%%%%%%%%%%

\begin{frame}{The Delta Method: Asymptotic Distribution}

Combine the Taylor expansion with $\hat{\bm{\beta}} \overset{a}{\sim} N(\bm{\beta},\ \operatorname{Asy.Var}(\hat{\bm{\beta}}))$:
\[
\bm{f}(\hat{\bm{\beta}}) \;\overset{a}{\sim}\; N\!\Bigl(\bm{f}(\bm{\beta}),\; \bm{\Gamma}\,[\operatorname{Asy.Var}(\hat{\bm{\beta}})]\,\bm{\Gamma}'\Bigr).
\]

\pause

\medskip

\textbf{Estimator.} Plug in $\bm{C}(\hat{\bm{\beta}})$ and the estimated variance of $\hat{\bm{\beta}}$:
\[
\widehat{\operatorname{Asy.Var}}\bigl(\bm{f}(\hat{\bm{\beta}})\bigr) \;=\; \bm{C}(\hat{\bm{\beta}})\,\bigl[s^2(\bm{X}'\bm{X})^{-1}\bigr]\,\bm{C}(\hat{\bm{\beta}})'.
\]

\pause

\medskip

\textit{HW4 Q1(b) is the scalar case: $g(x) = e^x$ gives $[g'(\mu)]^2\sigma^2 = \sigma^2 e^{2\mu}$.}

\end{frame}

%%%%%%%%%%%%%%%%%%%%%%%%%%%%%%%%%%%%%%%%%%%%%%%%%%%%%%%%%%%%

\begin{frame}{Long-Run Effects with a Lagged DV (Greene 4.4)}

Dynamic model: $y_t = \beta_0 + \beta_1 x_t + \gamma\, y_{t-1} + \epsilon_t$. Raise $x_t$ by 1 and trace forward (holding the system at its long-run mean otherwise):
\begin{align*}
y_t^* &= \bar y + \beta_1, &
y_{t+1}^* &= \bar y + \gamma\beta_1, &
y_{t+2}^* &= \bar y + \gamma^2\beta_1, &
\dots
\end{align*}
The cumulative long-run impact is a geometric series:
\[
\beta_1 + \gamma\beta_1 + \gamma^2\beta_1 + \cdots = \frac{\beta_1}{1-\gamma}.
\]
\textit{Equivalent shortcut:} solve $\bar y = \beta_0 + \beta_1 \bar x + \gamma \bar y$ at long-run equilibrium.

\end{frame}

%%%%%%%%%%%%%%%%%%%%%%%%%%%%%%%%%%%%%%%%%%%%%%%%%%%%%%%%%%%%

\begin{frame}{Application: Long-Run Gasoline Elasticities}

Log-log gasoline demand with lagged consumption,
\[
\log(G/\text{pop})_t = \beta_1 + \beta_2 \log P_{G,t} + \beta_3 \log(Y/\text{pop})_t + \beta_4 \log(G/\text{pop})_{t-1} + e_t.
\]

\begin{itemize}
\item Short-run price elasticity $= \beta_2$; short-run income elasticity $= \beta_3$.
\item By the lagged-DV iteration argument, long-run elasticities are
\[
\frac{\beta_2}{1-\beta_4} \quad\text{(price)}, \qquad \frac{\beta_3}{1-\beta_4} \quad\text{(income)}.
\]
\item Both are nonlinear in $\bm{\beta}$ --- next slide gets their standard errors via the delta method.
\end{itemize}

\end{frame}

%%%%%%%%%%%%%%%%%%%%%%%%%%%%%%%%%%%%%%%%%%%%%%%%%%%%%%%%%%%%

\begin{frame}{Delta Method for Long-Run Elasticities: Partials}

Stack the two long-run elasticities as a vector-valued function of $\bm{\beta} = (\beta_1, \beta_2, \beta_3, \beta_4)'$:
\[
\bm{f}(\bm{\beta}) \;=\; \begin{pmatrix} f_1(\bm{\beta}) \\ f_2(\bm{\beta}) \end{pmatrix}
\;=\; \begin{pmatrix} \beta_2/(1-\beta_4) \\ \beta_3/(1-\beta_4) \end{pmatrix}.
\]

\pause

\medskip

\textbf{Partials of $f_1 = \beta_2/(1-\beta_4)$:}
\begin{itemize}\setlength\itemsep{2pt}
\item $\partial f_1 / \partial \beta_1 = 0$ --- intercept does not enter.
\item $\partial f_1 / \partial \beta_2 = 1/(1-\beta_4)$ --- $f_1$ is linear in $\beta_2$ with that slope.
\item $\partial f_1 / \partial \beta_3 = 0$ --- $\beta_3$ does not enter.
\item $\partial f_1 / \partial \beta_4 = \beta_2/(1-\beta_4)^2$ --- chain rule on $1/(1-\beta_4)$.
\end{itemize}

\pause

\medskip

\textbf{Partials of $f_2 = \beta_3/(1-\beta_4)$.} Same pattern with the roles of $\beta_2$ and $\beta_3$ swapped: only $\partial f_2/\partial\beta_3$ and $\partial f_2/\partial\beta_4$ are nonzero.

\end{frame}

%%%%%%%%%%%%%%%%%%%%%%%%%%%%%%%%%%%%%%%%%%%%%%%%%%%%%%%%%%%%

\begin{frame}{Delta Method for Long-Run Elasticities: The Sandwich}

Stack the partials into the $2 \times 4$ Jacobian (one row per function, one column per coefficient):
\[
\bm{C}(\bm{\beta}) \;=\; \frac{\partial \bm f}{\partial \bm{\beta}'} \;=\;
\begin{bmatrix}
0 & \tfrac{1}{1-\beta_4} & 0 & \tfrac{\beta_2}{(1-\beta_4)^2} \\[2pt]
0 & 0 & \tfrac{1}{1-\beta_4} & \tfrac{\beta_3}{(1-\beta_4)^2}
\end{bmatrix}.
\]

\pause

\medskip

Plug in $\hat{\bm{\beta}}$ and sandwich with the estimated variance of $\hat{\bm{\beta}}$:
\[
\widehat{\Var}(\hat{\bm f}) \;=\; \bm{C}(\hat{\bm{\beta}})\,\hat{\bm{V}}_{\hat{\beta}}\,\bm{C}(\hat{\bm{\beta}})'.
\]

\pause

\medskip

Result is $2 \times 2$: diagonal entries are the variances of the long-run price and income elasticities, off-diagonal is their covariance.

\end{frame}

%%%%%%%%%%%%%%%%%%%%%%%%%%%%%%%%%%%%%%%%%%%%%%%%%%%%%%%%%%%%

\begin{frame}[fragile]{Gas Consumption: Implementation}

\begin{lstlisting}
fm <- lm(logG ~ logGASP + log(INCOME) + log(PNC)
              + log(PUC) + logGlag, data = USGasG)

# By hand
bs <- coef(fm); V <- vcov(fm)
g <- c(0, 1/(1-bs[6]), 0, 0, 0, bs[2]/(1-bs[6])^2)
se <- sqrt(t(g) %*% V %*% g)

# Canned: car::deltaMethod(fm, "logGASP/(1-logGlag)")
#   Estimate -0.411, SE 0.152, 95% CI [-0.710, -0.113]
# Stata: nlcom (_b[x2] / (1 - _b[x6]))
\end{lstlisting}

Long-run price elasticity $\approx -0.41 \pm 0.30$; income elasticity $\approx 0.97 \pm 0.32$ --- both larger in magnitude than short-run.

\end{frame}

%%%%%%%%%%%%%%%%%%%%%%%%%%%%%%%%%%%%%%%%%%%%%%%%%%%%%%%%%%%%

\begin{frame}{Delta Method for Probit: Average Marginal Effects}

Recall from Lecture 9: the Probit average marginal effect is a \textbf{nonlinear function} of $\hat{\bm{\beta}}$:
\[
\widehat{\text{AME}}_j = \frac{1}{n}\sum_{i=1}^n \phi(\bm{X}_i'\hat{\bm{\beta}})\,\hat{\beta}_j
\]

\pause

Apply the delta method. Define $h(\bm{\beta}) = \text{AME}_j(\bm{\beta})$. Then:
\[
\operatorname{Asy.Var}(\widehat{\text{AME}}_j) = \nabla h(\hat{\bm{\beta}})' \cdot \widehat{\operatorname{Var}}(\hat{\bm{\beta}}) \cdot \nabla h(\hat{\bm{\beta}})
\]

\pause

where $\nabla h$ accounts for both the direct effect ($\hat{\beta}_j$) and the indirect effect through $\phi(\bm{X}_i'\hat{\bm{\beta}})$.

\pause

\begin{tcolorbox}[colback=blue!5, colframe=blue!50]
The same delta-method machinery handles AMEs, long-run elasticities, and any other smooth transformation of $\hat{\bm{\theta}}$.
\end{tcolorbox}

\end{frame}

%%%%%%%%%%%%%%%%%%%%%%%%%%%%%%%%%%%%%%%%%%%%%%%%%%%%%%%%%%%%
\section{Consistency of OLS}
%%%%%%%%%%%%%%%%%%%%%%%%%%%%%%%%%%%%%%%%%%%%%%%%%%%%%%%%%%%%

\begin{frame}{Consistency of OLS: The Question}

\textbf{What we just assumed.} The delta method required $\hat{\bm{\beta}}$ to be close to $\bm{\beta}$ for large $n$. Is it?

\pause

\medskip

\textbf{What we are about to show.} Under the assumptions on the next slide,
\[
\hat{\bm{\beta}} \;\overset{p}{\to}\; \bm{\beta},
\]
i.e.\ $\Pr(\|\hat{\bm{\beta}} - \bm{\beta}\| > \epsilon) \to 0$ for every $\epsilon > 0$.

\pause

\medskip

\textbf{Tools}, all from Lecture 10:
\begin{itemize}\setlength\itemsep{2pt}
\item \textbf{WLLN}: sample moments converge in probability to population moments.
\item \textbf{CMT}: continuous functions of convergent sequences converge.
\end{itemize}
The next few slides apply these to $\hat{\bm{\beta}}$, written as a function of sample moments.

\end{frame}

%%%%%%%%%%%%%%%%%%%%%%%%%%%%%%%%%%%%%%%%%%%%%%%%%%%%%%%%%%%%

\begin{frame}{Target of Estimation}
\begin{itemize}
\item We are after the projection coefficient $\beta=(\mathbb{E}[\bm{x}\bm{x}'])^{-1}\mathbb{E}[\bm{x}Y]$ from the projection model $Y=\bm{x}'\beta+e$
\item We assume $(Y_i,\ \bm{x}_i)$ are i.i.d.
\item And define $\bm{Q}_{XX}=\mathbb{E}[\bm{x}\bm{x}']$, assume is positive definite.
\item We also need moments to exist, including 4th moments for asymptotic normality.
\end{itemize}
\end{frame}

%%%%%%%%%%%%%%%%%%%%%%%%%%%%%%%%%%%%%%%%%%%%%%%%%%%%%%%%%%%%

\begin{frame}{Proof Outline: Two Routes to Consistency}
\begin{itemize}
\item \textbf{Route 1} (Method of Moments):
\begin{itemize}
\item WLLN: sample moments converge to population moments
\item CMT: continuous functions of convergent sequences converge
\item OLS is a continuous function of sample moments
\end{itemize}
\pause
\item \textbf{Route 2} (Algebraic):
\begin{itemize}
\item Write $\hat{\bm{\beta}} - \bm{\beta}$ as a product of sample moments
\item Apply WLLN to each piece, CMT/Slutsky to combine
\end{itemize}
\end{itemize}

\pause

Both routes use the same tools from Lecture 10. Route 1 generalizes to nonlinear estimators (MLE, GMM); Route 2 gives an explicit formula for the sampling error.

\end{frame}

%%%%%%%%%%%%%%%%%%%%%%%%%%%%%%%%%%%%%%%%%%%%%%%%%%%%%%%%%%%%

\begin{frame}{OLS is a function of sample moments}
\begin{itemize}
\item $\hat{\bm{\beta}}=\left(\frac{1}{n}\sum_{i=1}^n\bm{x}_i\bm{x}'_i\right)^{-1}\left(\frac{1}{n}\sum_{i=1}^n\bm{x}_iY_i\right)=\hat{\bm{Q}}_{XX}^{-1}\hat{\bm{Q}}_{XY}$ \pause
\item $\hat{\bm{Q}}_{XX}$, $\hat{\bm{Q}}_{XY}$ are sample moments
\end{itemize}
\end{frame}

%%%%%%%%%%%%%%%%%%%%%%%%%%%%%%%%%%%%%%%%%%%%%%%%%%%%%%%%%%%%

\begin{frame}{Apply WLLN}
\begin{itemize}
\item If $(\bm{x}_i,\ Y_i)$ are iid, then $\bm{x}_i\bm{x}_i'$, $\bm{x}_iY_i'$ are iid.\pause
\item Recall from Lecture 10: WLLN says that if $\theta_i$ are iid with finite variance, $\frac{1}{n} \sum_{i=1}^n \theta_i \overset{p}{\to} \mathbb{E}[\theta]$\pause
\item $\hat{\bm{Q}}_{XX} \overset{p}{\to} \bm{Q}_{XX}$, \quad $\hat{\bm{Q}}_{XY} \overset{p}{\to} \bm{Q}_{XY}$
\end{itemize}
\end{frame}

%%%%%%%%%%%%%%%%%%%%%%%%%%%%%%%%%%%%%%%%%%%%%%%%%%%%%%%%%%%%

\begin{frame}{Continuous Mapping Theorem (CMT)}
\begin{itemize}
\item Recall from Lecture 10: if $Z_n\rightarrow_p c$ and $g(u)$ is continuous at $c$, then $g(Z_n)\rightarrow_p g(c)$.\pause
\item CMT tells us that plims survive continuous functions.\pause
\item $\hat{\bm{\beta}}=g(\hat{\bm{Q}}_{XX}, \hat{\bm{Q}}_{XY})$ is a continuous function (matrix inversion is continuous at nonsingular matrices).\pause
\item $\hat{\bm{\beta}}=\hat{\bm{Q}}_{XX}^{-1}\hat{\bm{Q}}_{XY}\rightarrow_p \bm{Q}_{XX}^{-1}\bm{Q}_{XY}=\bm{\beta}$
\end{itemize}
\end{frame}

%%%%%%%%%%%%%%%%%%%%%%%%%%%%%%%%%%%%%%%%%%%%%%%%%%%%%%%%%%%%

\begin{frame}{Proof of Consistency of OLS: Route 2, Step 1}
\begin{align*}
\hat{\bm{\beta}}&=(\bm{X'X})^{-1}\bm{X'y} \tag{Definition}\\ \pause
&=(\bm{X'X})^{-1}\bm{X}'(\bm{X}\bm{\beta}+\bm{e}) \tag{Assumption 1}\\  \pause
&=\bm{\beta}+(\bm{X'X})^{-1}\bm{X'e} \tag{Inverse prop.}\\  \pause
&=(\bm{\beta}+(1/n)n(\bm{X'X})^{-1}\bm{X'e}) \\ \pause
&=(\bm{\beta}+(\frac{1}{n}(\frac{1}{n})^{-1}(\bm{X'X})^{-1}\bm{X'e})) \tag{trick}\\ \pause
\hat{\bm{\beta}}&=\bm{\beta}+(\frac{1}{n}\bm{X'X})^{-1}) ( \frac{1}{n}\bm{X'e})
\end{align*}
 \end{frame}

%%%%%%%%%%%%%%%%%%%%%%%%%%%%%%%%%%%%%%%%%%%%%%%%%%%%%%%%%%%%

\begin{frame}{Proof of Consistency of OLS: Route 2, Step 2}
\begin{align*}
\hat{\bm{\beta}}&=\bm{\beta}+(\frac{1}{n}\bm{X'X})^{-1}) ( \frac{1}{n}\bm{X'e}) \\ \pause
\text{plim}\hat{\bm{\beta}}-\bm{\beta}&=((\text{plim}\frac{1}{n}\bm{X'X})^{-1}) \text{plim} \frac{1}{n}\bm{X'e} \tag{Slutsky}\\ \pause
&=\bm{Q}_{XX}^{-1} \mathbb{E}[Xe]\\
&=\bm{Q}_{XX}^{-1} 0
\end{align*}
The least squares estimator $\hat{\bm{\beta}}$ is \emph{consistent} for the projection coefficient $\bm{\beta}$.

\pause

\begin{tcolorbox}[colback=blue!5, colframe=blue!50]
\textbf{Key condition:} Consistency requires $\mathbb{E}[\bm{x}e] = 0$. When this fails (endogeneity), OLS is inconsistent --- motivating instrumental variables (Lecture 13).
\end{tcolorbox}

 \end{frame}

%%%%%%%%%%%%%%%%%%%%%%%%%%%%%%%%%%%%%%%%%%%%%%%%%%%%%%%%%%%%
\section{Asymptotic Normality}
%%%%%%%%%%%%%%%%%%%%%%%%%%%%%%%%%%%%%%%%%%%%%%%%%%%%%%%%%%%%

\begin{frame}{Asymptotic Normality: Algebraic Setup}
Start from the sampling error decomposition:
\begin{align*}
\hat{\bm{\beta}} - \bm{\beta}
&= \left( \frac{1}{n} \sum_{i=1}^n X_i X_i' \right)^{-1}
    \left( \frac{1}{n} \sum_{i=1}^n X_i e_i \right)
    \tag{1}
\end{align*}

\pause

Multiply both sides by $\sqrt{n}$:
\begin{align*}
\sqrt{n}(\hat{\bm{\beta}} - \bm{\beta})
&= \left( \frac{1}{n} \sum_{i=1}^n X_i X_i' \right)^{-1}
    \left( \frac{1}{\sqrt{n}} \sum_{i=1}^n X_i e_i \right)
    \tag{2} \\
&= \hat{\bm{Q}}_{XX}^{-1}
    \left( \frac{1}{\sqrt{n}} \sum_{i=1}^n X_i e_i \right)
    \tag{3}
\end{align*}

\pause

Now we need the distribution of $\frac{1}{\sqrt{n}} \sum X_i e_i$. This is a sum of mean-zero iid random vectors --- exactly the setting of the \textbf{multivariate CLT} from Lecture 10.

\end{frame}

%%%%%%%%%%%%%%%%%%%%%%%%%%%%%%%%%%%%%%%%%%%%%%%%%%%%%%%%%%%%

\begin{frame}{Asymptotic Normality: Applying CLT and Slutsky}
\begin{align*}
\sqrt{n}(\hat{\bm{\beta}} - \bm{\beta})
&= \underbrace{\hat{\bm{Q}}_{XX}^{-1}}_{\overset{p}{\to}\, \bm{Q}_{XX}^{-1}}
    \cdot \underbrace{\frac{1}{\sqrt{n}} \sum_{i=1}^n X_i e_i}_{\overset{d}{\to}\, N(0, \Omega)}
\end{align*}
where $\Omega = \mathbb{E}[X_i X_i' e_i^2]$.

\pause

By Slutsky's theorem ($\overset{p}{\to} \times \overset{d}{\to}$):
\begin{align*}
\sqrt{n}(\hat{\bm{\beta}} - \bm{\beta})
\xrightarrow{d}\; N\!\left( 0,\, \bm{Q}_{XX}^{-1} \Omega\, \bm{Q}_{XX}^{-1} \right)
    \equiv N(0,\, \bm{V}_{\beta})
\end{align*}

\pause

\begin{tcolorbox}[colback=blue!5, colframe=blue!50]
Same decomposition as Probit MLE in Lecture 9: CLT for the numerator, WLLN/CMT for the denominator, Slutsky to combine.
\end{tcolorbox}

\end{frame}

%%%%%%%%%%%%%%%%%%%%%%%%%%%%%%%%%%%%%%%%%%%%%%%%%%%%%%%%%%%%

\begin{frame}{Covariance of coefficient estimators under homoskedasticity}
\begin{itemize}
\item Assume $Y=\beta_1X_1+\beta_2X_2+e$ with $E[ee']=\sigma^2\bm{I}$.
\item Suppose $E[X_1]=0$, $E[X_2]=0$, $Var[X_1]=1$, $Var[X_2]=1$, $corr[X_1, X_2]=\rho$
$$\bm{V}_\beta^0=\sigma^2\bm{Q}_{XX}^{-1}=\sigma^2\begin{bmatrix}1&\rho\\ \rho &1 \end{bmatrix}^{-1}=\frac{\sigma^2}{1-\rho^2}\begin{bmatrix}1&-\rho\\ -\rho &1 \end{bmatrix}$$
\item If $X_1$ and $X_2$ are positively correlated, $\hat{\beta}_1$ and $\hat{\beta}_2$ are negatively correlated.
\end{itemize}
\end{frame}

%%%%%%%%%%%%%%%%%%%%%%%%%%%%%%%%%%%%%%%%%%%%%%%%%%%%%%%%%%%%

\begin{frame}{Heteroskedastic Covariance Matrix Estimation}
\begin{itemize}
\item The asymptotic covariance matrix of $\sqrt{n}(\hat{\bm{\beta}}-\bm{\beta})$ is $\bm{Q}_{XX}^{-1} \Omega \bm{Q}_{XX}^{-1}$.
\item This is the \textbf{sandwich form} from Lecture 6 --- now with its asymptotic justification.
\end{itemize}

\pause

\textbf{Estimator:}
$$\hat{\Omega}=\frac{1}{n}\sum_{i=1}^n \bm{x}_i\bm{x}_i'\hat{e}_i^2$$
$$\hat{\bm{Q}}_{XX}=\frac{1}{n}\sum_{i=1}^n\bm{x}_i\bm{x}_i'$$
$$\hat{\bm{V}}_\beta^{HC0}=\hat{\bm{Q}}_{XX}^{-1} \hat{\Omega} \hat{\bm{Q}}_{XX}^{-1}$$

\pause

This is the HC0 estimator. Recall from Lecture 6: HC1, HC2, HC3 apply finite-sample corrections. All are asymptotically equivalent.
\end{frame}

%%%%%%%%%%%%%%%%%%%%%%%%%%%%%%%%%%%%%%%%%%%%%%%%%%%%%%%%%%%%

\begin{frame}{Consistency of $\hat{\Omega}$}
Add and subtract population errors:
\begin{align*}
\hat{\Omega}
&=\underbrace{\frac{1}{n}\sum_{i=1}^n \bm{x}_i\bm{x}_i'e_i^2}_{\overset{p}{\to}\,\Omega \text{ (WLLN)}}+\underbrace{\frac{1}{n}\sum_{i=1}^n \bm{x}_i\bm{x}_i'(\hat{e}_i^2-e_i^2)}_{\overset{p}{\to}\, 0}
\end{align*}

\pause

The remainder vanishes: $\hat{e}_i - e_i = \bm{X}_i'(\bm{\beta} - \hat{\bm{\beta}}) \overset{p}{\to} 0$ by consistency of $\hat{\bm{\beta}}$.

\pause

Therefore $\hat{\Omega}\overset{p}{\to}\Omega$, and by CMT:
$$\hat{\bm{V}}_\beta^{HC0}=\hat{\bm{Q}}_{XX}^{-1} \hat{\Omega}\, \hat{\bm{Q}}_{XX}^{-1} \overset{p}{\to} \bm{Q}_{XX}^{-1}\Omega\,\bm{Q}_{XX}^{-1} = \bm{V}_\beta$$

\end{frame}

%%%%%%%%%%%%%%%%%%%%%%%%%%%%%%%%%%%%%%%%%%%%%%%%%%%%%%%%%%%%

\begin{frame}{The Sandwich Pattern}

The same ``bread--meat--bread'' structure appears throughout this course:

\medskip

\renewcommand{\arraystretch}{1.5}
\begin{tabular}{lccc}
\toprule
\textbf{Model} & \textbf{Bread} & \textbf{Meat} ($\Omega$) & \textbf{Simplification} \\
\midrule
OLS & $Q_{XX}^{-1}$ & $\mathbb{E}[XX'e^2]$ & $\sigma^2 Q_{XX}^{-1}$ if homosk. \\
Probit MLE & $\mathcal{I}^{-1}$ & $\mathbb{E}[s_i s_i']$ & $\mathcal{I}^{-1}$ if correct spec. \\
GMM & $(G'WG)^{-1}G'W$ & $\mathbb{E}[g_i g_i']$ & --- \\
\bottomrule
\end{tabular}
\renewcommand{\arraystretch}{1.0}

\pause

\bigskip

\begin{itemize}
    \item Under ``ideal'' conditions (homoskedasticity / correct specification), the sandwich simplifies.
    \item Under misspecification, we need the full sandwich --- \textbf{robust SEs}.
    \item Same \texttt{sandwich::vcovHC()} in R handles all cases.
\end{itemize}

\end{frame}

%%%%%%%%%%%%%%%%%%%%%%%%%%%%%%%%%%%%%%%%%%%%%%%%%%%%%%%%%%%%

\begin{frame}{Variances: Exact vs Asymptotic}

\renewcommand{\arraystretch}{1.4}
\begin{tabular}{lll}
\toprule
\textbf{Object} & \textbf{Formula} & \textbf{When to use} \\
\midrule
Exact variance & $(\bm{X}'\bm{X})^{-1}(\bm{X}'\bm{D}\bm{X})(\bm{X}'\bm{X})^{-1}$ & Fixed $X$, any $n$ \\
\quad (homosk.) & $\sigma^2(\bm{X}'\bm{X})^{-1}$ & $+ E[ee']=\sigma^2 I$ \\
\midrule
Asy. variance & $\bm{Q}_{XX}^{-1}\Omega\bm{Q}_{XX}^{-1}$ & Random $X$, $n$ large \\
\quad (homosk.) & $\sigma^2 \bm{Q}_{XX}^{-1}$ & $+ E[e^2|X]=\sigma^2$ \\
\bottomrule
\end{tabular}
\renewcommand{\arraystretch}{1.0}

\pause

\bigskip

\begin{itemize}
    \item The exact variance conditions on $X$; the asymptotic variance averages over $X$.
    \item For large $n$: $\frac{1}{n}\bm{X'X} \approx Q_{XX}$, so they agree.
    \item In practice, we always estimate using the sandwich (or its homoskedastic special case).
\end{itemize}

\end{frame}

%%%%%%%%%%%%%%%%%%%%%%%%%%%%%%%%%%%%%%%%%%%%%%%%%%%%%%%%%%%%
\section{Asymptotic Properties: Wald Test}
\setcounter{subsection}{1}
%%%%%%%%%%%%%%%%%%%%%%%%%%%%%%%%%%%%%%%%%%%%%%%%%%%%%%%%%%%%

\begin{frame}{From F Tests to Wald Tests}
\begin{itemize}
\item We saw we can calculate an F test by forming a matrix of restrictions.
\item The F distribution result depended on the homoskedasticity and normality of the errors.
\item In large samples, we can instead rely on the asymptotic behavior of random variables and calculate a \textbf{Wald statistic}.
\end{itemize}

\pause

\begin{itemize}
\item Many test statistics are reflavored Wald distances. Given a null that $E[\bm{q}]=\bm{\theta}$:
\begin{align*}
W &= (\bm{q}-\bm{\theta})'[\widehat{Var}(\bm{q}-\bm{\theta})]^{-1}(\bm{q}-\bm{\theta})
\end{align*}
\item $W \sim\chi^2_J$ if $\bm{q}$ is Normal and $Var(\bm{q}-\bm{\theta})=\bm{\Sigma}$
\item $W \overset{a}{\sim}\chi^2_J$ with CLT on $\bm{q}$ and consistency for $\hat{\bm{\Sigma}}$
\end{itemize}
\end{frame}

%%%%%%%%%%%%%%%%%%%%%%%%%%%%%%%%%%%%%%%%%%%%%%%%%%%%%%%%%%%%

\begin{frame}{Mahalanobis Distance: Intuition}
\begin{itemize}
\item Suppose $\bm{x} \sim \mathcal{N}(\bm{\mu}, \bm{\Sigma})$. We want to measure how unusual an observation $\bm{x}$ is.

\item In the scalar case:
\[
D = \frac{x - \mu}{\sigma} \sim \mathcal{N}(0,1) \quad \Rightarrow \quad D^2 \sim \chi^2_1
\]

\item In the vector case: transform the data to make it standard normal:
\[
\bm{z} = \bm{\Sigma}^{-1/2} (\bm{x} - \bm{\mu}) \sim \mathcal{N}(0, \bm{I})
\]

\item Then the squared norm of $\bm{z}$ gives:
\[
\|\bm{z}\|^2 = (\bm{x} - \bm{\mu})' \bm{\Sigma}^{-1} (\bm{x} - \bm{\mu}) = D^2
\]

\item This is the Mahalanobis distance: it tells us how far $\bm{x}$ is from $\bm{\mu}$ in standardized units, accounting for correlations and scale.
\end{itemize}
\end{frame}

%%%%%%%%%%%%%%%%%%%%%%%%%%%%%%%%%%%%%%%%%%%%%%%%%%%%%%%%%%%%

\begin{frame}{Wald Statistic as a Distance}
\begin{itemize}
\item The Wald test generalizes the standardized $t$-statistic to higher dimensions:
\[
\text{Univariate: } \frac{\hat{\beta} - \beta_0}{\text{SE}(\hat{\beta})} \sim \mathcal{N}(0,1)
\quad \Rightarrow \quad \left( \frac{\hat{\beta} - \beta_0}{\text{SE}(\hat{\beta})} \right)^2 \sim \chi^2_1
\]
\item In multivariate form:
\[
W = (\bm{q} - \bm{\theta})' \hat{\bm{\Sigma}}^{-1} (\bm{q} - \bm{\theta})
\]
\item If $\bm{q}$ is consistent and asymptotically normal:
\[
\sqrt{n}(\bm{q} - \bm{\theta}) \overset{d}{\to} \mathcal{N}(0, \bm{\Sigma}) \Rightarrow W \overset{a}{\sim} \chi^2_J
\]
\item Interpretation: Wald tests ask whether the observed $\bm{q}$ is ``far'' from $\bm{\theta}$, accounting for sampling variability.
\end{itemize}
\end{frame}

%%%%%%%%%%%%%%%%%%%%%%%%%%%%%%%%%%%%%%%%%%%%%%%%%%%%%%%%%%%%

\begin{frame}{Matrix form of Restrictions}
\begin{itemize}
\item Suppose we have $k$ parameters $\bm{\beta}$ and $q$ restrictions with $q\leq (k)$. (Note $q=J$ in Greene, here k includes the intercept)
\item Let $\bm{R}$ be a $q\times(k)$ matrix.
\item We can express our null hypothesis as $H_0:\ \bm{R\beta}=\bm{r}$.
\item For example in the prior example where $\beta_2=0$ and $\beta_3=\beta_4$\\
$$\bm{R}=\begin{bmatrix} 0& 0&1& 0&0\\ 0&0& 0& 1&-1\end{bmatrix},\quad \bm{r}=\begin{bmatrix} 0\\0\end{bmatrix}$$
$$\begin{bmatrix} 0&0& 1& 0&0\\0& 0& 0& 1&-1\end{bmatrix} \begin{bmatrix} \beta_0\\\beta_1\\\beta_2\\\beta_3\\\beta_4 \end{bmatrix}=\begin{bmatrix} 0\\0\end{bmatrix}$$
\end{itemize}
\end{frame}

%%%%%%%%%%%%%%%%%%%%%%%%%%%%%%%%%%%%%%%%%%%%%%%%%%%%%%%%%%%%

\begin{frame}{Wald Test for Linear Restrictions}
$$W=(\bm{R}\hat{\bm{\beta}}-\bm{r})'[\bm{R}\hat{\bm{V}}_{\hat{\beta}}\bm{R}']^{-1}(\bm{R}\hat{\bm{\beta}}-\bm{r})$$
$$W \overset{a}{\sim} \chi^2_q$$
\begin{itemize}
\item If we use $\hat{\bm{V}}_{\hat{\beta}} = s^2(\bm{X}'\bm{X})^{-1}$ (homoskedasticity + normality), then $W/q$ is exactly the $F$ statistic distributed $F_{q,n-k}$.
\item Without normality, the CLT tells us $W/q$ is approximately $F_{q,n-k}$.
\item Without normality or homoskedasticity, use $\hat{\bm{V}}_{\hat{\beta}}^{HC}$ (sandwich) and $W \overset{a}{\sim} \chi^2_q$.
\end{itemize}

\pause

\begin{tcolorbox}[colback=blue!5, colframe=blue!50]
The Wald test nests $F$. With robust SEs, only consistency + asymptotic normality required. \textit{Single coef: $t^2 = F$ (HW4 Q3d).}
\end{tcolorbox}

\end{frame}

%%%%%%%%%%%%%%%%%%%%%%%%%%%%%%%%%%%%%%%%%%%%%%%%%%%%%%%%%%%%

\begin{frame}[fragile]{Wald Tests in R}
  \begin{lstlisting}
library(sandwich); library(lmtest); library(car)

mod0 <- lm(prestige ~ income + education + women,
            data = Prestige)

# Wald test with robust SEs (chi-squared)
linearHypothesis(mod0,
    c("education = 0", "women = 0"),
    vcov = vcovHC, test = "Chisq")

# Equivalent: waldtest with nested models
mod1 <- lm(prestige ~ income, data = Prestige)
waldtest(mod1, mod0, vcov = vcovHC, test = "Chisq")
\end{lstlisting}

\pause

Both test $H_0: \beta_{ed} = \beta_{women} = 0$ using the sandwich covariance.

\end{frame}

%%%%%%%%%%%%%%%%%%%%%%%%%%%%%%%%%%%%%%%%%%%%%%%%%%%%%%%%%%%%
\section{Delta Method Applied to OLS}
\setcounter{subsection}{1}
%%%%%%%%%%%%%%%%%%%%%%%%%%%%%%%%%%%%%%%%%%%%%%%%%%%%%%%%%%%%

\begin{frame}{Applying the Delta Method to OLS}

Earlier we introduced the delta method in general: if $\hat{\bm{\theta}}$ is asymptotically normal, then any smooth function $h(\hat{\bm{\theta}})$ is also asymptotically normal.

\pause

\bigskip

Now we have the specific ingredients for OLS:
\begin{itemize}
  \item $\hat{\bm{\beta}}$ is asymptotically normal (proved in the previous sections)
  \item $\hat{\bm{V}}_{\hat{\beta}}$ is a consistent estimator of its covariance (sandwich or homoskedastic)
  \item For any differentiable $\theta = r(\bm{\beta})$, the gradient $\bm{R} = \frac{\partial r}{\partial \bm{\beta}}\big|_{\hat{\bm{\beta}}}$ gives:
\end{itemize}

\begin{tcolorbox}
$$s(\hat{\theta})=\sqrt{\bm{R}'\hat{\bm{V}}_{\hat{\beta}}\bm{R}}$$
This works for \textbf{any} differentiable function of $\hat{\bm{\beta}}$ --- linear or nonlinear.
\end{tcolorbox}

\pause

Let's work through three examples of increasing complexity.

\end{frame}

%%%%%%%%%%%%%%%%%%%%%%%%%%%%%%%%%%%%%%%%%%%%%%%%%%%%%%%%%%%%

\begin{frame}{Running Example: Log Wage Regression}
\begin{itemize}
\item $\log(wage)=\beta_1 \cdot education+\beta_2 \cdot experience+\beta_3 \cdot experience^2/100+\beta_4+e$.
\item We can calculate:
\begin{itemize}
\item Percentage return to education:
$$\theta_1=100\beta_1$$
\item Percentage return to experience for those with 10 years of experience:
$$\theta_2=100\beta_2+20\beta_3$$
\item Experience level which maximizes expected log wages:
$$\theta_3=-50\beta_2/\beta_3$$
\end{itemize}
\end{itemize}

\pause

Each is a function $\theta = r(\beta)$. The first two are linear; the third is nonlinear. All three use the same delta method formula.

\end{frame}

%%%%%%%%%%%%%%%%%%%%%%%%%%%%%%%%%%%%%%%%%%%%%%%%%%%%%%%%%%%%

\begin{frame}{Example: Return to Education (Linear)}
\begin{itemize}
\item Percentage return to education:
$$\theta_1=100\beta_1$$
\item $\bm{R}=\begin{pmatrix}100\\ 0\\ 0\\ 0\end{pmatrix}$
\item Asymptotic standard error:
\[
s(\hat{\theta}_1) = \sqrt{\bm{R}' \hat{\bm{V}}_{\hat{\beta}} \bm{R}} = 100 \cdot \sqrt{\hat{\text{Var}}(\hat{\beta}_1)}
\]
\end{itemize}

This is just rescaling --- the SE scales linearly.

\end{frame}

%%%%%%%%%%%%%%%%%%%%%%%%%%%%%%%%%%%%%%%%%%%%%%%%%%%%%%%%%%%%

\begin{frame}{Example: Return to Experience at 10 Years (Linear)}
\begin{itemize}
\item Percentage return to experience at 10 years:
\[
\theta_2 = 100\beta_2 + 20\beta_3
\]
\item Gradient of $\theta_2$ with respect to $\bm{\beta}$:
\[
\bm{R} =
\begin{pmatrix}
0 \\
100 \\
20 \\
0
\end{pmatrix}
\]
\item Asymptotic standard error:
\[
s(\hat{\theta}_2) = \sqrt{\bm{R}' \hat{\bm{V}}_{\hat{\beta}} \bm{R}}
\]
\end{itemize}

This is a linear combination of two coefficients --- the SE depends on their covariance.

\end{frame}

%%%%%%%%%%%%%%%%%%%%%%%%%%%%%%%%%%%%%%%%%%%%%%%%%%%%%%%%%%%%

\begin{frame}{Example: Maximum Return to Experience (Nonlinear)}
\begin{itemize}
\item Experience level which maximizes expected log wages:
\[
\theta_3 = -50\frac{\beta_2}{\beta_3}
\]
\item Gradient of $\theta_3$ with respect to $\bm{\beta}$:
\[
\bm{R} =
\begin{pmatrix}
0 \\
-50/\beta_3 \\
50\beta_2/\beta_3^2 \\
0
\end{pmatrix}
\]
\item Asymptotic standard error:
\[
s(\hat{\theta}_3) = \sqrt{\bm{R}' \hat{\bm{V}}_{\hat{\beta}} \bm{R}}
\]
\end{itemize}

This is nonlinear --- $\bm{R}$ itself depends on $\bm{\beta}$, estimated at $\hat{\bm{\beta}}$.

\end{frame}

%%%%%%%%%%%%%%%%%%%%%%%%%%%%%%%%%%%%%%%%%%%%%%%%%%%%%%%%%%%%

\begin{frame}{The z-statistic}

In asymptotic inference the $t$-statistic is called the $z$-statistic:
\[
T(\theta) = \frac{\hat{\theta} - \theta}{s(\hat{\theta})}
= \frac{\sqrt n(\hat{\theta} - \theta)}{\sqrt{\hat{V}_\theta}}
\xrightarrow{d} N(0,1).
\]
Under $H_0: \theta = \theta_0$, reject at level $\alpha$ if $|T| > z_{\alpha/2}$ (e.g.\ $z_{0.025} = 1.96$).

\end{frame}

%%%%%%%%%%%%%%%%%%%%%%%%%%%%%%%%%%%%%%%%%%%%%%%%%%%%%%%%%%%%

\begin{frame}{Regression Intervals}

For the CEF $m(\bm{x}) = \bm{x}'\bm{\beta}$, the target parameter $\theta = r(\bm{\beta}) = \bm{x}'\bm{\beta}$ has gradient $\bm{R} = \bm{x}$, so the delta method gives
\[
s(\hat{\theta}) = \sqrt{\bm{x}' \hat{\bm{V}}_{\hat{\beta}}\, \bm{x}}.
\]
A $95\%$ CI for $\mathbb{E}[Y \mid \bm{X} = \bm{x}]$ is
\[
\bm{x}'\hat{\bm{\beta}} \;\pm\; 1.96 \cdot \sqrt{\bm{x}' \hat{\bm{V}}_{\hat{\beta}}\, \bm{x}}.
\]

\end{frame}

%%%%%%%%%%%%%%%%%%%%%%%%%%%%%%%%%%%%%%%%%%%%%%%%%%%%%%%%%%%%

\begin{frame}{Example: Prestige $\sim$ Education}

$\hat{\beta}_{ed}=5.4$, $\hat{\beta}_{cons}=-10.7$,  $\hat{\bm{V}}_{\hat{\beta}}=\begin{bmatrix} 13.52&-1.18\\-1.18&.11\end{bmatrix}$\\
at 10 years of education:
$$\bm{x}' \hat{\bm{\beta}} \pm 1.96 \cdot \sqrt{\bm{x}' \hat{\bm{V}}_{\hat{\beta}} \bm{x}}$$
$$(1\ 10) \begin{pmatrix}-10.7\\ 5.4\end{pmatrix} \pm 1.96 \cdot \sqrt{(1\ 10) \begin{bmatrix} 13.52&-1.18\\-1.18&.11\end{bmatrix}\begin{pmatrix}1\\ 10\end{pmatrix}}$$
$$43.3 \pm 1.96 \cdot \sqrt{0.87} = 43.3 \pm 1.83$$
compare to 22 years of education: $\;107 \pm 1.96 \cdot \sqrt{14.8} = 107 \pm 7.5$
\end{frame}

%%%%%%%%%%%%%%%%%%%%%%%%%%%%%%%%%%%%%%%%%%%%%%%%%%%%%%%%%%%%

\begin{frame}[fragile]{R: Sandwich SEs and Delta Method}

\begin{lstlisting}
# Model-based vs sandwich SEs
coeftest(mod0)
coeftest(mod0, vcov = vcovHC(mod0, type = "HC1"))

# CI for CEF at a point
x0 <- c(1, 5000, 12, 30)
se <- sqrt(t(x0) %*% vcovHC(mod0) %*% x0)
cat(sum(x0 * coef(mod0)), "+/-", 1.96 * se)

# Delta method for a nonlinear function
car::deltaMethod(mod0, "income / education")
\end{lstlisting}

\end{frame}

%%%%%%%%%%%%%%%%%%%%%%%%%%%%%%%%%%%%%%%%%%%%%%%%%%%%%%%%%%%%
\section{Application: Ridge Regression}
%%%%%%%%%%%%%%%%%%%%%%%%%%%%%%%%%%%%%%%%%%%%%%%%%%%%%%%%%%%%

\begin{frame}{Application: Ridge Regression}
\begin{itemize}
\item Suppose $Y=\bm{X}'\bm{\beta}+e$ with $\mathbb{E}[\bm{X}e]=\bm{0}$
$$\hat{\bm{\beta}}_R=\left(\sum_{i=1}^n \bm{X}_i\bm{X}_i'+\lambda \bm{I}_k\right)^{-1}\left(\sum_{i=1}^n \bm{X}_iY_i\right)$$
\item Where $\lambda>0$ is a constant, added to the diagonals of the covariance matrix.
\item This is biased, shrinks $\hat{\bm{\beta}}$.
\item Arises from solving $\min_{\bm{\beta}} (\bm{y}-\bm{X}\bm{\beta})'(\bm{y}-\bm{X}\bm{\beta})+\lambda(\bm{\beta}'\bm{\beta}-c)$
\end{itemize}
\end{frame}

%%%%%%%%%%%%%%%%%%%%%%%%%%%%%%%%%%%%%%%%%%%%%%%%%%%%%%%%%%%%

\begin{frame}{Application: Ridge Regression}
\begin{center}
\includegraphics[width=2.5in, height=0.75\textheight, keepaspectratio]{ridge.png}
\end{center}
\end{frame}

%%%%%%%%%%%%%%%%%%%%%%%%%%%%%%%%%%%%%%%%%%%%%%%%%%%%%%%%%%%%

\begin{frame}{Ridge is Consistent (via CMT)}
\begin{align*}
\text{plim}\ \hat{\bm{\beta}}_R
&= \text{plim} \left( \frac{1}{n} \sum_{i=1}^n \bm{X}_i \bm{X}_i' + \frac{\lambda}{n} \bm{I}_k \right)^{-1}
\left( \frac{1}{n} \sum_{i=1}^n \bm{X}_i Y_i \right) \\[4pt]
&= \left( \mathbb{E}[\bm{X}_i \bm{X}_i'] + \underbrace{\text{plim}\,\frac{\lambda}{n}}_{\to\, 0}\, \bm{I}_k \right)^{-1}
\mathbb{E}[\bm{X}_i Y_i] \\[4pt]
&= \left( \mathbb{E}[\bm{X}_i \bm{X}_i'] \right)^{-1} \mathbb{E}[\bm{X}_i Y_i] = \bm{\beta}
\end{align*}

\pause
\begin{itemize}
  \item As $n \to \infty$, Ridge converges to the same plim as OLS.
  \item Biased in finite samples, consistent asymptotically.
\end{itemize}
\end{frame}

%%%%%%%%%%%%%%%%%%%%%%%%%%%%%%%%%%%%%%%%%%%%%%%%%%%%%%%%%%%%
\section{Summary}
%%%%%%%%%%%%%%%%%%%%%%%%%%%%%%%%%%%%%%%%%%%%%%%%%%%%%%%%%%%%

\begin{frame}{What We Proved Today}

The OLS estimator $\hat{\bm{\beta}}=(\bm{X}'\bm{X})^{-1}\bm{X}'y$ is:

\medskip

\renewcommand{\arraystretch}{1.4}
\begin{tabular}{lll}
\toprule
\textbf{Property} & \textbf{Tool (from Lectures 10--11)} & \textbf{Result} \\
\midrule
Consistent & WLLN + CMT & $\text{plim}\ \hat{\bm{\beta}} = \bm{\beta}$ \\
Asy. normal & Multivariate CLT + Slutsky & $\sqrt{n}(\hat{\bm{\beta}} - \bm{\beta}) \overset{d}{\to} N(\bm{0}, \bm{V}_\beta)$ \\
Testable & Wald statistic & $W \overset{a}{\sim} \chi^2_q$ \\
Robust & Sandwich estimator & $\hat{\bm{V}}_\beta^{HC} \overset{p}{\to} \bm{V}_\beta$ \\
\bottomrule
\end{tabular}
\renewcommand{\arraystretch}{1.0}

\pause

\bigskip

These results hold \textbf{without normality} --- only requiring iid data with finite 4th moments and $\mathbb{E}[\bm{x}e] = 0$.

\end{frame}

%%%%%%%%%%%%%%%%%%%%%%%%%%%%%%%%%%%%%%%%%%%%%%%%%%%%%%%%%%%%

\begin{frame}{Beyond OLS: the OLS Argument as a Template}

\textbf{Recap what we just did.} In the normality step we wrote
\[
\sqrt n(\hat{\bm{\beta}} - \bm{\beta}) \;=\; \frac{1}{\sqrt n}\sum_{i=1}^n \underbrace{\bm{Q}_{XX}^{-1}\,\bm{x}_i\,e_i}_{\psi(Z_i)} \;+\; o_p(1),
\]
a sample mean of mean-zero random vectors plus a negligible remainder. CLT $+$ Slutsky then delivered asymptotic normality.

\pause

\medskip

The whole argument turned on two things:
\begin{enumerate}\setlength\itemsep{2pt}
\item Expressing $\sqrt n(\hat{\theta} - \theta)$ as a sample mean of an iid, mean-zero random quantity.
\item Controlling the leftover as $o_p(1)$.
\end{enumerate}
Nothing about the argument was specific to OLS. Next slide: the general version.

\end{frame}

%%%%%%%%%%%%%%%%%%%%%%%%%%%%%%%%%%%%%%%%%%%%%%%%%%%%%%%%%%%%

\begin{frame}{Influence Functions (Preview)}

\textbf{The general template.} Many estimators $\hat{\theta}$ admit the same representation:
\[
\sqrt n(\hat{\theta} - \theta) \;=\; \frac{1}{\sqrt n}\sum_{i=1}^n \psi(Z_i) \;+\; o_p(1), \qquad \mathbb{E}[\psi(Z)] = 0.
\]
The function $\psi$ is the estimator's \textbf{influence function}: it measures how much a single observation $Z_i$ moves $\hat{\theta}$ around $\theta$.

\pause

\medskip

CLT on the sample mean, Slutsky to absorb the $o_p(1)$:
\[
\sqrt n(\hat{\theta} - \theta) \;\overset{d}{\to}\; N\bigl(0,\ \mathrm{Var}(\psi)\bigr).
\]

\pause

\medskip

\emph{HW4 Q4--Q9 builds this template for the AIPW estimator of the average treatment effect.}

\end{frame}

%%%%%%%%%%%%%%%%%%%%%%%%%%%%%%%%%%%%%%%%%%%%%%%%%%%%%%%%%%%%

\begin{frame}{Looking Ahead}

\textbf{What happens when $\mathbb{E}[\bm{x}e] \neq 0$?}

\pause

\begin{itemize}
    \item OLS is \textbf{inconsistent}: $\text{plim}\,\hat{\bm{\beta}} = \bm{\beta} + \bm{Q}_{XX}^{-1}\mathbb{E}[\bm{X}e] \neq \bm{\beta}$
    \item No amount of data can fix this --- the bias does not vanish.
\end{itemize}

\pause

\bigskip

\textbf{Lectures 13--14 (IV and 2SLS):} Find instruments $Z$ with $\mathbb{E}[Ze] = 0$ and $\mathbb{E}[ZX'] \neq 0$.
\begin{itemize}
    \item Replace the OLS moment condition $\mathbb{E}[X(Y - X'\beta)] = 0$ with $\mathbb{E}[Z(Y - X'\beta)] = 0$
    \item Same asymptotic machinery: WLLN, CLT, Slutsky, sandwich --- applied to a different estimating equation
    \item This is another instance of the GMM framework: $\mathbb{E}[g(W_i, \theta_0)] = 0$
\end{itemize}

\end{frame}


\end{document}
